% paper90-copilot-advisors.tex — The Advisor Architecture:
%   Domain-Specific Copilot Guidance in JuGeo.
%   Formalises five domain-specific advisors (heap, scope, import,
%   callable, contract) that translate judgment-geometry analysis into
%   actionable Copilot suggestions, proves trust-ceiling and soundness
%   theorems, and evaluates the architecture experimentally.
%   Paper 90 of the Judgment Geometry series.
% Compile: pdflatex paper90-copilot-advisors

\documentclass[11pt]{article}
\usepackage{lmodern}
% jugeo-common.tex — shared preamble for all Judgment Geometry papers
% Include via: % jugeo-common.tex — shared preamble for all Judgment Geometry papers
% Include via: % jugeo-common.tex — shared preamble for all Judgment Geometry papers
% Include via: \input{jugeo-common}

% ── Packages ────────────────────────────────────────────────────────────
\usepackage{amsmath,amssymb,amsthm,mathtools}
\usepackage{tikz-cd}
\usepackage{listings}
\usepackage{xcolor}
\usepackage{xspace}
\usepackage{booktabs}
\usepackage{multirow}
\usepackage{enumitem}
\usepackage[expansion=false]{microtype}
\usepackage{hyperref}
\usepackage{cleveref}
% \usepackage{thmtools}  % disabled: not available in this TeX installation
\usepackage{float}
\usepackage{caption}
\usepackage{subcaption}
\usepackage{tabularx}
\usepackage{stmaryrd}       % \llbracket, \rrbracket
\usepackage{bbm}            % \mathbbm{1}

% ── Page geometry ───────────────────────────────────────────────────────
\usepackage[margin=1in]{geometry}

% ── Theorem environments ────────────────────────────────────────────────
\theoremstyle{plain}
\newtheorem{theorem}{Theorem}[section]
\newtheorem{lemma}[theorem]{Lemma}
\newtheorem{proposition}[theorem]{Proposition}
\newtheorem{corollary}[theorem]{Corollary}

\theoremstyle{definition}
\newtheorem{definition}[theorem]{Definition}
\newtheorem{example}[theorem]{Example}
\newtheorem{construction}[theorem]{Construction}

\theoremstyle{remark}
\newtheorem{remark}[theorem]{Remark}
\newtheorem{notation}[theorem]{Notation}

% ── Python listings style ──────────────────────────────────────────────
\definecolor{jugeoBlue}{HTML}{2563EB}
\definecolor{jugeoGreen}{HTML}{059669}
\definecolor{jugeoRed}{HTML}{DC2626}
\definecolor{jugeoPurple}{HTML}{7C3AED}
\definecolor{jugeoGray}{HTML}{6B7280}
\definecolor{jugeoBg}{HTML}{F8FAFC}

\lstdefinestyle{jugeo-python}{
  language=Python,
  basicstyle=\ttfamily\small,
  keywordstyle=\color{jugeoBlue}\bfseries,
  stringstyle=\color{jugeoGreen},
  commentstyle=\color{jugeoGray}\itshape,
  numberstyle=\tiny\color{jugeoGray},
  backgroundcolor=\color{jugeoBg},
  numbers=left,
  numbersep=6pt,
  frame=single,
  framerule=0.4pt,
  rulecolor=\color{jugeoGray!40},
  breaklines=true,
  breakatwhitespace=true,
  showstringspaces=false,
  tabsize=4,
  xleftmargin=1.5em,
  framexleftmargin=1.5em,
  morekeywords={dataclass,frozen,field,Enum,IntEnum,auto,Any,
                Mapping,Sequence,Optional,match,case},
  literate={→}{{$\to$}}1 {←}{{$\leftarrow$}}1
           {≤}{{$\leq$}}1 {≥}{{$\geq$}}1 {≠}{{$\neq$}}1
           {∈}{{$\in$}}1 {∀}{{$\forall$}}1 {∃}{{$\exists$}}1
           {⊕}{{$\oplus$}}1 {⊖}{{$\ominus$}}1,
  escapeinside={(*@}{@*)},
}
\lstset{style=jugeo-python}

\lstdefinestyle{jugeo-output}{
  basicstyle=\ttfamily\small,
  backgroundcolor=\color{jugeoBg},
  frame=single,
  framerule=0.4pt,
  rulecolor=\color{jugeoGray!40},
  breaklines=true,
  showstringspaces=false,
  xleftmargin=1.5em,
  framexleftmargin=0.5em,
  numbers=none,
}

% ── JuGeo-specific macros ──────────────────────────────────────────────

% Judgment 8-tuple
\newcommand{\J}{\mathcal{J}}
\newcommand{\Jud}[8]{(#1,\,#2,\,#3,\,#4,\,#5,\,#6,\,#7,\,#8)}
\newcommand{\JudShort}[1]{\J_{#1}}

% Components
\newcommand{\coord}{\mathsf{c}}            % coordinate
\newcommand{\prop}{\varphi}                 % proposition
\newcommand{\carrier}{\mathcal{A}}          % carrier type
\newcommand{\evid}{\mathcal{E}}             % evidence bundle
\newcommand{\oblig}{\mathcal{O}}            % obligations
\newcommand{\obstr}{\mathcal{B}}            % obstructions
\newcommand{\trust}{\mathcal{T}}            % trust annotation
\newcommand{\prov}{\Pi}                     % provenance

% Trust algebra
\newcommand{\Talg}{\mathfrak{T}}            % trust ordered algebra
\newcommand{\Eadm}{\mathcal{E}_{\mathrm{adm}}}
\newcommand{\tleq}{\preceq}                 % trust ordering
\newcommand{\tjoin}{\oplus}                 % conservative join
\newcommand{\tatten}{\ominus}               % attenuation
\newcommand{\tpromote}{\uparrow_\pi}        % promotion
\newcommand{\tdemote}{\downarrow_\chi}       % demotion

% Trust tiers
\newcommand{\Tcontra}{\ensuremath{\mathsf{CONTRADICTED}}}
\newcommand{\Tunver}{\ensuremath{\mathsf{UNVERIFIED}}}
\newcommand{\Tcopilot}{\ensuremath{\mathsf{COPILOT}}}
\newcommand{\Toracle}{\ensuremath{\mathsf{ORACLE}}}
\newcommand{\Truntime}{\ensuremath{\mathsf{RUNTIME}}}
\newcommand{\Tsolver}{\ensuremath{\mathsf{SOLVER}}}
\newcommand{\Tproof}{\ensuremath{\mathsf{PROOF}}}

% Site and geometry
\newcommand{\Site}{\mathcal{S}}             % semantic site
\newcommand{\Cov}{\mathrm{Cov}}            % covering family
\newcommand{\Ob}{\mathrm{Ob}}              % objects
\newcommand{\Mor}{\mathrm{Mor}}            % morphisms
\newcommand{\res}{\mathrm{res}}            % restriction
\newcommand{\Cech}{\check{C}}              % Čech complex
\newcommand{\Hcoh}[1]{H^{#1}}             % cohomology group

% Descent
\newcommand{\descent}{\mathrm{desc}}
\newcommand{\glue}{\mathrm{glue}}
\newcommand{\overlap}[2]{#1 \cap #2}

% Semantic moves
\newcommand{\VERIFY}{\textsc{Verify}}
\newcommand{\CONSTRUCT}{\textsc{Construct}}
\newcommand{\NORMALIZE}{\textsc{Normalize}}
\newcommand{\GLUE}{\textsc{Glue}}
\newcommand{\RESTRICT}{\textsc{Restrict}}
\newcommand{\TRANSPORT}{\textsc{Transport}}
\newcommand{\REPAIR}{\textsc{Repair}}
\newcommand{\PROMOTE}{\textsc{Promote}}
\newcommand{\CHALLENGE}{\textsc{Challenge}}
\newcommand{\ESCALATE}{\textsc{Escalate}}

% Fragments
\newcommand{\QFLIA}{\texttt{QF\_LIA}}
\newcommand{\QFLRA}{\texttt{QF\_LRA}}
\newcommand{\QFBV}{\texttt{QF\_BV}}
\newcommand{\QFUF}{\texttt{QF\_UF}}

% Misc
\newcommand{\Python}{\textsc{Python}}
\newcommand{\JuGeo}{\textsc{JuGeo}}
\newcommand{\LEAN}{\textsc{Lean}\,4}
\newcommand{\Fstar}{F$^{\star}$}
\newcommand{\Coq}{\textsc{Coq}}
\newcommand{\Dafny}{\textsc{Dafny}}

% Common shortcuts
\newcommand{\ie}{i.e.\@\xspace}
\newcommand{\eg}{e.g.\@\xspace}
\newcommand{\cf}{cf.\@\xspace}
\newcommand{\wrt}{w.r.t.\@\xspace}

% Evaluation shortcuts
\newcommand{\numcases}{300}
\newcommand{\numspec}{100}
\newcommand{\numeq}{100}
\newcommand{\numbug}{100}
\newcommand{\medtime}{6.5\,ms}
\newcommand{\totaltime}{1.949\,s}


% ── Packages ────────────────────────────────────────────────────────────
\usepackage{amsmath,amssymb,amsthm,mathtools}
\usepackage{tikz-cd}
\usepackage{listings}
\usepackage{xcolor}
\usepackage{xspace}
\usepackage{booktabs}
\usepackage{multirow}
\usepackage{enumitem}
\usepackage[expansion=false]{microtype}
\usepackage{hyperref}
\usepackage{cleveref}
% \usepackage{thmtools}  % disabled: not available in this TeX installation
\usepackage{float}
\usepackage{caption}
\usepackage{subcaption}
\usepackage{tabularx}
\usepackage{stmaryrd}       % \llbracket, \rrbracket
\usepackage{bbm}            % \mathbbm{1}

% ── Page geometry ───────────────────────────────────────────────────────
\usepackage[margin=1in]{geometry}

% ── Theorem environments ────────────────────────────────────────────────
\theoremstyle{plain}
\newtheorem{theorem}{Theorem}[section]
\newtheorem{lemma}[theorem]{Lemma}
\newtheorem{proposition}[theorem]{Proposition}
\newtheorem{corollary}[theorem]{Corollary}

\theoremstyle{definition}
\newtheorem{definition}[theorem]{Definition}
\newtheorem{example}[theorem]{Example}
\newtheorem{construction}[theorem]{Construction}

\theoremstyle{remark}
\newtheorem{remark}[theorem]{Remark}
\newtheorem{notation}[theorem]{Notation}

% ── Python listings style ──────────────────────────────────────────────
\definecolor{jugeoBlue}{HTML}{2563EB}
\definecolor{jugeoGreen}{HTML}{059669}
\definecolor{jugeoRed}{HTML}{DC2626}
\definecolor{jugeoPurple}{HTML}{7C3AED}
\definecolor{jugeoGray}{HTML}{6B7280}
\definecolor{jugeoBg}{HTML}{F8FAFC}

\lstdefinestyle{jugeo-python}{
  language=Python,
  basicstyle=\ttfamily\small,
  keywordstyle=\color{jugeoBlue}\bfseries,
  stringstyle=\color{jugeoGreen},
  commentstyle=\color{jugeoGray}\itshape,
  numberstyle=\tiny\color{jugeoGray},
  backgroundcolor=\color{jugeoBg},
  numbers=left,
  numbersep=6pt,
  frame=single,
  framerule=0.4pt,
  rulecolor=\color{jugeoGray!40},
  breaklines=true,
  breakatwhitespace=true,
  showstringspaces=false,
  tabsize=4,
  xleftmargin=1.5em,
  framexleftmargin=1.5em,
  morekeywords={dataclass,frozen,field,Enum,IntEnum,auto,Any,
                Mapping,Sequence,Optional,match,case},
  literate={→}{{$\to$}}1 {←}{{$\leftarrow$}}1
           {≤}{{$\leq$}}1 {≥}{{$\geq$}}1 {≠}{{$\neq$}}1
           {∈}{{$\in$}}1 {∀}{{$\forall$}}1 {∃}{{$\exists$}}1
           {⊕}{{$\oplus$}}1 {⊖}{{$\ominus$}}1,
  escapeinside={(*@}{@*)},
}
\lstset{style=jugeo-python}

\lstdefinestyle{jugeo-output}{
  basicstyle=\ttfamily\small,
  backgroundcolor=\color{jugeoBg},
  frame=single,
  framerule=0.4pt,
  rulecolor=\color{jugeoGray!40},
  breaklines=true,
  showstringspaces=false,
  xleftmargin=1.5em,
  framexleftmargin=0.5em,
  numbers=none,
}

% ── JuGeo-specific macros ──────────────────────────────────────────────

% Judgment 8-tuple
\newcommand{\J}{\mathcal{J}}
\newcommand{\Jud}[8]{(#1,\,#2,\,#3,\,#4,\,#5,\,#6,\,#7,\,#8)}
\newcommand{\JudShort}[1]{\J_{#1}}

% Components
\newcommand{\coord}{\mathsf{c}}            % coordinate
\newcommand{\prop}{\varphi}                 % proposition
\newcommand{\carrier}{\mathcal{A}}          % carrier type
\newcommand{\evid}{\mathcal{E}}             % evidence bundle
\newcommand{\oblig}{\mathcal{O}}            % obligations
\newcommand{\obstr}{\mathcal{B}}            % obstructions
\newcommand{\trust}{\mathcal{T}}            % trust annotation
\newcommand{\prov}{\Pi}                     % provenance

% Trust algebra
\newcommand{\Talg}{\mathfrak{T}}            % trust ordered algebra
\newcommand{\Eadm}{\mathcal{E}_{\mathrm{adm}}}
\newcommand{\tleq}{\preceq}                 % trust ordering
\newcommand{\tjoin}{\oplus}                 % conservative join
\newcommand{\tatten}{\ominus}               % attenuation
\newcommand{\tpromote}{\uparrow_\pi}        % promotion
\newcommand{\tdemote}{\downarrow_\chi}       % demotion

% Trust tiers
\newcommand{\Tcontra}{\ensuremath{\mathsf{CONTRADICTED}}}
\newcommand{\Tunver}{\ensuremath{\mathsf{UNVERIFIED}}}
\newcommand{\Tcopilot}{\ensuremath{\mathsf{COPILOT}}}
\newcommand{\Toracle}{\ensuremath{\mathsf{ORACLE}}}
\newcommand{\Truntime}{\ensuremath{\mathsf{RUNTIME}}}
\newcommand{\Tsolver}{\ensuremath{\mathsf{SOLVER}}}
\newcommand{\Tproof}{\ensuremath{\mathsf{PROOF}}}

% Site and geometry
\newcommand{\Site}{\mathcal{S}}             % semantic site
\newcommand{\Cov}{\mathrm{Cov}}            % covering family
\newcommand{\Ob}{\mathrm{Ob}}              % objects
\newcommand{\Mor}{\mathrm{Mor}}            % morphisms
\newcommand{\res}{\mathrm{res}}            % restriction
\newcommand{\Cech}{\check{C}}              % Čech complex
\newcommand{\Hcoh}[1]{H^{#1}}             % cohomology group

% Descent
\newcommand{\descent}{\mathrm{desc}}
\newcommand{\glue}{\mathrm{glue}}
\newcommand{\overlap}[2]{#1 \cap #2}

% Semantic moves
\newcommand{\VERIFY}{\textsc{Verify}}
\newcommand{\CONSTRUCT}{\textsc{Construct}}
\newcommand{\NORMALIZE}{\textsc{Normalize}}
\newcommand{\GLUE}{\textsc{Glue}}
\newcommand{\RESTRICT}{\textsc{Restrict}}
\newcommand{\TRANSPORT}{\textsc{Transport}}
\newcommand{\REPAIR}{\textsc{Repair}}
\newcommand{\PROMOTE}{\textsc{Promote}}
\newcommand{\CHALLENGE}{\textsc{Challenge}}
\newcommand{\ESCALATE}{\textsc{Escalate}}

% Fragments
\newcommand{\QFLIA}{\texttt{QF\_LIA}}
\newcommand{\QFLRA}{\texttt{QF\_LRA}}
\newcommand{\QFBV}{\texttt{QF\_BV}}
\newcommand{\QFUF}{\texttt{QF\_UF}}

% Misc
\newcommand{\Python}{\textsc{Python}}
\newcommand{\JuGeo}{\textsc{JuGeo}}
\newcommand{\LEAN}{\textsc{Lean}\,4}
\newcommand{\Fstar}{F$^{\star}$}
\newcommand{\Coq}{\textsc{Coq}}
\newcommand{\Dafny}{\textsc{Dafny}}

% Common shortcuts
\newcommand{\ie}{i.e.\@\xspace}
\newcommand{\eg}{e.g.\@\xspace}
\newcommand{\cf}{cf.\@\xspace}
\newcommand{\wrt}{w.r.t.\@\xspace}

% Evaluation shortcuts
\newcommand{\numcases}{300}
\newcommand{\numspec}{100}
\newcommand{\numeq}{100}
\newcommand{\numbug}{100}
\newcommand{\medtime}{6.5\,ms}
\newcommand{\totaltime}{1.949\,s}


% ── Packages ────────────────────────────────────────────────────────────
\usepackage{amsmath,amssymb,amsthm,mathtools}
\usepackage{tikz-cd}
\usepackage{listings}
\usepackage{xcolor}
\usepackage{xspace}
\usepackage{booktabs}
\usepackage{multirow}
\usepackage{enumitem}
\usepackage[expansion=false]{microtype}
\usepackage{hyperref}
\usepackage{cleveref}
% \usepackage{thmtools}  % disabled: not available in this TeX installation
\usepackage{float}
\usepackage{caption}
\usepackage{subcaption}
\usepackage{tabularx}
\usepackage{stmaryrd}       % \llbracket, \rrbracket
\usepackage{bbm}            % \mathbbm{1}

% ── Page geometry ───────────────────────────────────────────────────────
\usepackage[margin=1in]{geometry}

% ── Theorem environments ────────────────────────────────────────────────
\theoremstyle{plain}
\newtheorem{theorem}{Theorem}[section]
\newtheorem{lemma}[theorem]{Lemma}
\newtheorem{proposition}[theorem]{Proposition}
\newtheorem{corollary}[theorem]{Corollary}

\theoremstyle{definition}
\newtheorem{definition}[theorem]{Definition}
\newtheorem{example}[theorem]{Example}
\newtheorem{construction}[theorem]{Construction}

\theoremstyle{remark}
\newtheorem{remark}[theorem]{Remark}
\newtheorem{notation}[theorem]{Notation}

% ── Python listings style ──────────────────────────────────────────────
\definecolor{jugeoBlue}{HTML}{2563EB}
\definecolor{jugeoGreen}{HTML}{059669}
\definecolor{jugeoRed}{HTML}{DC2626}
\definecolor{jugeoPurple}{HTML}{7C3AED}
\definecolor{jugeoGray}{HTML}{6B7280}
\definecolor{jugeoBg}{HTML}{F8FAFC}

\lstdefinestyle{jugeo-python}{
  language=Python,
  basicstyle=\ttfamily\small,
  keywordstyle=\color{jugeoBlue}\bfseries,
  stringstyle=\color{jugeoGreen},
  commentstyle=\color{jugeoGray}\itshape,
  numberstyle=\tiny\color{jugeoGray},
  backgroundcolor=\color{jugeoBg},
  numbers=left,
  numbersep=6pt,
  frame=single,
  framerule=0.4pt,
  rulecolor=\color{jugeoGray!40},
  breaklines=true,
  breakatwhitespace=true,
  showstringspaces=false,
  tabsize=4,
  xleftmargin=1.5em,
  framexleftmargin=1.5em,
  morekeywords={dataclass,frozen,field,Enum,IntEnum,auto,Any,
                Mapping,Sequence,Optional,match,case},
  literate={→}{{$\to$}}1 {←}{{$\leftarrow$}}1
           {≤}{{$\leq$}}1 {≥}{{$\geq$}}1 {≠}{{$\neq$}}1
           {∈}{{$\in$}}1 {∀}{{$\forall$}}1 {∃}{{$\exists$}}1
           {⊕}{{$\oplus$}}1 {⊖}{{$\ominus$}}1,
  escapeinside={(*@}{@*)},
}
\lstset{style=jugeo-python}

\lstdefinestyle{jugeo-output}{
  basicstyle=\ttfamily\small,
  backgroundcolor=\color{jugeoBg},
  frame=single,
  framerule=0.4pt,
  rulecolor=\color{jugeoGray!40},
  breaklines=true,
  showstringspaces=false,
  xleftmargin=1.5em,
  framexleftmargin=0.5em,
  numbers=none,
}

% ── JuGeo-specific macros ──────────────────────────────────────────────

% Judgment 8-tuple
\newcommand{\J}{\mathcal{J}}
\newcommand{\Jud}[8]{(#1,\,#2,\,#3,\,#4,\,#5,\,#6,\,#7,\,#8)}
\newcommand{\JudShort}[1]{\J_{#1}}

% Components
\newcommand{\coord}{\mathsf{c}}            % coordinate
\newcommand{\prop}{\varphi}                 % proposition
\newcommand{\carrier}{\mathcal{A}}          % carrier type
\newcommand{\evid}{\mathcal{E}}             % evidence bundle
\newcommand{\oblig}{\mathcal{O}}            % obligations
\newcommand{\obstr}{\mathcal{B}}            % obstructions
\newcommand{\trust}{\mathcal{T}}            % trust annotation
\newcommand{\prov}{\Pi}                     % provenance

% Trust algebra
\newcommand{\Talg}{\mathfrak{T}}            % trust ordered algebra
\newcommand{\Eadm}{\mathcal{E}_{\mathrm{adm}}}
\newcommand{\tleq}{\preceq}                 % trust ordering
\newcommand{\tjoin}{\oplus}                 % conservative join
\newcommand{\tatten}{\ominus}               % attenuation
\newcommand{\tpromote}{\uparrow_\pi}        % promotion
\newcommand{\tdemote}{\downarrow_\chi}       % demotion

% Trust tiers
\newcommand{\Tcontra}{\ensuremath{\mathsf{CONTRADICTED}}}
\newcommand{\Tunver}{\ensuremath{\mathsf{UNVERIFIED}}}
\newcommand{\Tcopilot}{\ensuremath{\mathsf{COPILOT}}}
\newcommand{\Toracle}{\ensuremath{\mathsf{ORACLE}}}
\newcommand{\Truntime}{\ensuremath{\mathsf{RUNTIME}}}
\newcommand{\Tsolver}{\ensuremath{\mathsf{SOLVER}}}
\newcommand{\Tproof}{\ensuremath{\mathsf{PROOF}}}

% Site and geometry
\newcommand{\Site}{\mathcal{S}}             % semantic site
\newcommand{\Cov}{\mathrm{Cov}}            % covering family
\newcommand{\Ob}{\mathrm{Ob}}              % objects
\newcommand{\Mor}{\mathrm{Mor}}            % morphisms
\newcommand{\res}{\mathrm{res}}            % restriction
\newcommand{\Cech}{\check{C}}              % Čech complex
\newcommand{\Hcoh}[1]{H^{#1}}             % cohomology group

% Descent
\newcommand{\descent}{\mathrm{desc}}
\newcommand{\glue}{\mathrm{glue}}
\newcommand{\overlap}[2]{#1 \cap #2}

% Semantic moves
\newcommand{\VERIFY}{\textsc{Verify}}
\newcommand{\CONSTRUCT}{\textsc{Construct}}
\newcommand{\NORMALIZE}{\textsc{Normalize}}
\newcommand{\GLUE}{\textsc{Glue}}
\newcommand{\RESTRICT}{\textsc{Restrict}}
\newcommand{\TRANSPORT}{\textsc{Transport}}
\newcommand{\REPAIR}{\textsc{Repair}}
\newcommand{\PROMOTE}{\textsc{Promote}}
\newcommand{\CHALLENGE}{\textsc{Challenge}}
\newcommand{\ESCALATE}{\textsc{Escalate}}

% Fragments
\newcommand{\QFLIA}{\texttt{QF\_LIA}}
\newcommand{\QFLRA}{\texttt{QF\_LRA}}
\newcommand{\QFBV}{\texttt{QF\_BV}}
\newcommand{\QFUF}{\texttt{QF\_UF}}

% Misc
\newcommand{\Python}{\textsc{Python}}
\newcommand{\JuGeo}{\textsc{JuGeo}}
\newcommand{\LEAN}{\textsc{Lean}\,4}
\newcommand{\Fstar}{F$^{\star}$}
\newcommand{\Coq}{\textsc{Coq}}
\newcommand{\Dafny}{\textsc{Dafny}}

% Common shortcuts
\newcommand{\ie}{i.e.\@\xspace}
\newcommand{\eg}{e.g.\@\xspace}
\newcommand{\cf}{cf.\@\xspace}
\newcommand{\wrt}{w.r.t.\@\xspace}

% Evaluation shortcuts
\newcommand{\numcases}{300}
\newcommand{\numspec}{100}
\newcommand{\numeq}{100}
\newcommand{\numbug}{100}
\newcommand{\medtime}{6.5\,ms}
\newcommand{\totaltime}{1.949\,s}

% data-paper90.tex — AUTO-GENERATED by exp90_copilot_advisors.py
% DO NOT EDIT — regenerate with: python3 experiments/exp90_copilot_advisors.py
% Generated from 12 programs across 5 advisor domains

\newcommand{\ppXCprogramCount}{12}
\newcommand{\ppXCprogramsOk}{12}
\newcommand{\ppXCverified}{12}
\newcommand{\ppXCverifiedPct}{100.0\%}

% ── Heap Advisor ────────────────────────────────────────────────────
\newcommand{\ppXCheapAdviceCount}{47}
\newcommand{\ppXCheapImmutSuggestions}{18}
\newcommand{\ppXCheapAliasingExplained}{12}
\newcommand{\ppXCheapMutationBugs}{9}
\newcommand{\ppXCheapCopySuggestions}{8}
\newcommand{\ppXCheapAcceptRate}{82.4\%}
\newcommand{\ppXCheapTimeMean}{0.34\,s}

% ── Scope Advisor ───────────────────────────────────────────────────
\newcommand{\ppXCscopeAdviceCount}{39}
\newcommand{\ppXCscopeRenameSuggestions}{14}
\newcommand{\ppXCscopeShadowDetections}{11}
\newcommand{\ppXCscopeRefactorSuggestions}{8}
\newcommand{\ppXCscopeAnnotations}{6}
\newcommand{\ppXCscopeAcceptRate}{79.5\%}
\newcommand{\ppXCscopeTimeMean}{0.28\,s}

% ── Import Advisor ──────────────────────────────────────────────────
\newcommand{\ppXCimportCycleAdvice}{7}
\newcommand{\ppXCimportStarAdvice}{11}
\newcommand{\ppXCimportDynamicAdvice}{5}
\newcommand{\ppXCimportAcceptRate}{85.7\%}
\newcommand{\ppXCimportTimeMean}{0.19\,s}

% ── Callable Advisor ────────────────────────────────────────────────
\newcommand{\ppXCcallableAdviceCount}{34}
\newcommand{\ppXCcallableTypeSuggestions}{15}
\newcommand{\ppXCcallableDescriptorExplained}{8}
\newcommand{\ppXCcallableBindingErrors}{6}
\newcommand{\ppXCcallableRefactorSuggestions}{5}
\newcommand{\ppXCcallableAcceptRate}{76.5\%}
\newcommand{\ppXCcallableTimeMean}{0.41\,s}

% ── Contract Advisor ────────────────────────────────────────────────
\newcommand{\ppXCcontractAdviceCount}{52}
\newcommand{\ppXCcontractMissingAnnotations}{22}
\newcommand{\ppXCcontractFixAnnotations}{17}
\newcommand{\ppXCcontractProposals}{13}
\newcommand{\ppXCcontractAcceptRate}{88.5\%}
\newcommand{\ppXCcontractTimeMean}{0.37\,s}

% ── Aggregate ───────────────────────────────────────────────────────
\newcommand{\ppXCtotalAdvice}{195}
\newcommand{\ppXCtotalAcceptRate}{82.6\%}
\newcommand{\ppXCtotalBugsFound}{26}
\newcommand{\ppXCcomposedAdvisors}{5}
\newcommand{\ppXCtimeMean}{0.32\,s}
\newcommand{\ppXCtimeTotal}{3.84\,s}
\newcommand{\ppXCtrustCeiling}{0.91}
\newcommand{\ppXCsoundnessRate}{97.4\%}


% ── Additional packages ────────────────────────────────────────────
\usepackage{array}
\usepackage{tikz}
\usetikzlibrary{arrows.meta,positioning,calc,fit,backgrounds}
\usepackage{algorithm}
\usepackage{algorithmic}

% ── Paper-specific macros ──────────────────────────────────────────
\newcommand{\Adv}{\mathcal{A}}
\newcommand{\AdvH}{\Adv_{\mathsf{H}}}
\newcommand{\AdvS}{\Adv_{\mathsf{S}}}
\newcommand{\AdvI}{\Adv_{\mathsf{I}}}
\newcommand{\AdvC}{\Adv_{\mathsf{C}}}
\newcommand{\AdvG}{\Adv_{\mathsf{G}}}
\newcommand{\AdvComp}{\Adv^{\otimes}}
\newcommand{\Domain}{\mathcal{D}}
\newcommand{\Trust}{\tau}
\newcommand{\TrustCeil}{\overline{\Trust}}
\newcommand{\Conf}{\kappa}
\newcommand{\Advice}{\mathsf{Advice}}
\newcommand{\Sound}{\mathsf{Sound}}
\newcommand{\Inv}{\mathcal{I}}
\newcommand{\Accept}{\alpha}
\newcommand{\Coverage}{\gamma}
\newcommand{\AdvisorOp}{\mathsf{advise}}
\newcommand{\ComposeOp}{\mathsf{compose}}
\newcommand{\HeapPart}{\Pi_{\mathsf{H}}}
\newcommand{\ScopeRes}{\Sigma_{\mathsf{S}}}
\newcommand{\ImportBundle}{\mathcal{B}_{\mathsf{I}}}
\newcommand{\CallSurf}{\mathcal{F}_{\mathsf{C}}}
\newcommand{\ContractRes}{\mathcal{R}_{\mathsf{G}}}
\newcommand{\DomainSet}{\{\mathsf{H},\mathsf{S},\mathsf{I},\mathsf{C},\mathsf{G}\}}

\title{\textbf{The Advisor Architecture:\\
Domain-Specific Copilot Guidance in JuGeo}}

\author{%
  JuGeo Research Group
}

\date{\today}

\begin{document}
\maketitle

\paragraph{Audit note.}
The current repository contains the underlying analysis domains and a
small results artifact, but not a stable exported five-advisor API.
Accordingly, any remaining interface sketches in this paper should be
read as design sketches rather than as importable JuGeo classes.

% ─────────────────────────────────────────────────────────────────────
\begin{abstract}
JuGeo contains several analysis subsystems whose outputs could be
surfaced as developer-facing guidance, but the repository does not
currently export the exact five-class advisor API described in the
original draft.  This audited version therefore treats the paper as a
design note grounded in existing subsystems rather than as a report on a
shipped interface.  We keep the discussion of heap, scope, import,
callable, and contract-oriented guidance as motivating domains, and we
retain the small JSON artifact in \texttt{results\_paper90.json} as an
auxiliary example, but we remove claims of reproduced benchmark
performance, theorem proofs, and a stable \texttt{CopilotAdvisor}
hierarchy.
\end{abstract}

% =====================================================================
\section{Introduction}\label{sec:intro}
% =====================================================================

Static analysis engines and type checkers produce rich diagnostic data,
but developers rarely consume raw analysis output.  Modern integrated
development environments address this gap with \emph{code actions},
\emph{quick fixes}, and \emph{Copilot suggestions}—actionable items
that translate an analysis finding into a concrete edit.  The
effectiveness of such guidance depends on three properties:

\begin{enumerate}
  \item \textbf{Trust.}  The developer must believe that the suggestion
    is correct.  A single false positive can erode trust for an entire
    tool category.
  \item \textbf{Relevance.}  The suggestion must address the developer's
    current concern, not a tangentially related issue.
  \item \textbf{Composability.}  When multiple analyses fire
    simultaneously, their suggestions should compose without conflict.
\end{enumerate}

JuGeo's judgment-geometry pipeline~\cite{paper00} already provides
categorical models of five Python runtime domains: heap
aliasing~\cite{paper32}, scope resolution~\cite{paper31}, import
graphs~\cite{paper49}, callable surfaces~\cite{paper61}, and generated
contracts~\cite{paper70}.  Each domain produces structured analysis
results—partitions, resolution records, import bundles, surface
descriptors, contract annotations—that are mathematically precise but
not immediately useful to a developer working in a Copilot-assisted
environment.

This paper therefore presents the \emph{advisor architecture} as a
proposed integration layer over existing JuGeo analyses rather than as a
completed library surface.  Each advisor sketch consumes a
domain-specific analysis result and emits a prioritised list of
$\Advice$ records, but the exact class hierarchy should be treated as
prospective.

\paragraph{Contributions.}
\begin{enumerate}
  \item We define the common advisor interface and show how each of
    the five domain-specific advisors specialises it
    (Section~\ref{sec:pattern}).
  \item We present $\AdvH$ and $\AdvS$ in detail, covering heap
    immutability suggestions, aliasing explanations, scope shadowing
    detection, and rename recommendations
    (Section~\ref{sec:heap-scope}).
  \item We present $\AdvI$ and $\AdvC$, covering import-cycle advice,
    star-import warnings, descriptor explanations, and binding-error
    detection (Section~\ref{sec:import-callable}).
  \item We present $\AdvG$, the contract advisor, covering missing
    annotation detection, fix proposals, and full type-annotation
    generation (Section~\ref{sec:contract}).
  \item We isolate the repository-backed trust and provenance constraints
    that any such advisor layer would need to obey
    (Section~\ref{sec:formal}).
  \item We retain the experiments section only as supporting methodology
    around the checked-in example results file
    (Section~\ref{sec:experiments}).
\end{enumerate}

\paragraph{Paper outline.}
Section~\ref{sec:pattern} introduces the advisor pattern.
Sections~\ref{sec:heap-scope}--\ref{sec:contract} detail each advisor.
Section~\ref{sec:formal} states and proves the formal properties.
Section~\ref{sec:experiments} presents experimental results.
Section~\ref{sec:related} surveys related work.
Section~\ref{sec:conclusion} concludes.

% =====================================================================
\section{The Advisor Pattern in JuGeo}\label{sec:pattern}
% =====================================================================

All five advisors share a common structural pattern.  In this section
we define the abstract interface and explain how domain specialisation
is achieved.

% ---------------------------------------------------------------------
\subsection{Common Interface}\label{sec:common-interface}
% ---------------------------------------------------------------------

An \emph{advisor} is a triple $\Adv = (\Domain, \TrustCeil, \AdvisorOp)$
where $\Domain$ is the analysis domain (one of $\DomainSet$),
$\TrustCeil \in [0,1]$ is the trust ceiling, and $\AdvisorOp$ is a
function that maps an analysis result to a list of $\Advice$ records.

\begin{definition}[Advice Record]\label{def:advice}
An advice record is a tuple $(\Domain, \Conf, \mathit{content},
\mathit{actionable})$ where $\Conf \in [0,\TrustCeil]$ is the
confidence, $\mathit{content}$ is a human-readable string, and
$\mathit{actionable} \in \{\top, \bot\}$ indicates whether the
advice maps to a concrete code edit.
\end{definition}

In Python, the common interface is realised by the following pattern
shared across all five advisor classes:

\begin{lstlisting}[style=jugeo-python, caption={Common advisor pattern (schematic).}, label={lst:common}]
class CopilotAdvisorBase:
    """Abstract base for all domain-specific Copilot advisors."""
    def __init__(self):
        self._advice_log: list[dict] = []
        self._enabled: bool = True

    def enable(self):
        self._enabled = True

    def disable(self):
        self._enabled = False

    def all_advice(self) -> list[dict]:
        return list(self._advice_log)

    def _record(self, kind: str, confidence: float, content: str,
                actionable: bool = True):
        if self._enabled:
            self._advice_log.append({
                "kind": kind,
                "confidence": min(confidence, self.TRUST_CEILING),
                "content": content,
                "actionable": actionable,
            })
\end{lstlisting}

The \texttt{TRUST\_CEILING} class attribute is set individually by each
advisor; for instance, $\AdvH$ uses $\TrustCeil = 0.90$ because heap
analysis may produce approximate alias partitions, while $\AdvG$ uses
$\TrustCeil = 0.95$ because generated contracts are derived from
concrete type inference.

\begin{remark}
The \texttt{\_record} helper enforces the trust ceiling invariant
at the point of advice creation, not at query time.  This ensures
that even if the analysis engine overestimates confidence, the
advice shown to the developer is clamped.
\end{remark}

The common interface also supports a \emph{disable/enable} toggle,
allowing individual advisors to be turned off during composition
without removing them from the advisor list.  This is important for
the composability theorem (Section~\ref{sec:composability}), which
quantifies only over enabled advisors.

Algorithm~\ref{alg:advise} formalises the common advice-generation
loop that all five advisors follow.

\begin{algorithm}[ht]
\caption{Common Advisor Loop}\label{alg:advise}
\begin{algorithmic}[1]
\REQUIRE Analysis result $R$, domain $\Domain$, trust ceiling $\TrustCeil$
\ENSURE Advice list $L$
\STATE $L \leftarrow []$
\FOR{each finding $f$ in $R$}
  \STATE $\Conf \leftarrow \mathsf{computeConfidence}(f)$
  \STATE $\Conf \leftarrow \min(\Conf, \TrustCeil)$
  \STATE $\mathit{content} \leftarrow \mathsf{explain}(f)$
  \STATE $\mathit{act} \leftarrow \mathsf{isActionable}(f)$
  \STATE $L.\mathsf{append}((\Domain, \Conf, \mathit{content}, \mathit{act}))$
\ENDFOR
\RETURN $L$
\end{algorithmic}
\end{algorithm}

The three helper functions—$\mathsf{computeConfidence}$,
$\mathsf{explain}$, and $\mathsf{isActionable}$—are the points
of domain specialisation.  Each advisor overrides them to reflect
its domain's semantics.

% ---------------------------------------------------------------------
\subsection{Domain Specialisation}\label{sec:specialisation}
% ---------------------------------------------------------------------

Each advisor subclass adds domain-specific methods.  The specialisation
is guided by two principles:

\begin{enumerate}
  \item \textbf{Analysis-driven.}  Every public method consumes a
    data structure produced by the corresponding analysis pass—a heap
    partition, a scope-resolution record, an import bundle, a callable
    surface, or a contract result.
  \item \textbf{Explanation-first.}  Before suggesting an edit, the
    advisor should be able to \emph{explain} the analysis finding in
    natural language.  This ensures that the developer can evaluate
    the suggestion's merit.
\end{enumerate}

The five domains and their advisor specialisations are summarised in
Table~\ref{tab:domains}.

\begin{table}[ht]
\centering
\caption{Advisor domains and their specialisations.}
\label{tab:domains}
\begin{tabular}{@{}llll@{}}
\toprule
\textbf{Domain} & \textbf{Class} & \textbf{Key Methods} & $\TrustCeil$ \\
\midrule
Heap ($\mathsf{H}$)
  & $\mathsf{CopilotHeapAdvisor}$
  & \texttt{suggest\_immutability}, \texttt{detect\_mutation\_bugs}
  & 0.90 \\
Scope ($\mathsf{S}$)
  & $\mathsf{CopilotScopeAdvisor}$
  & \texttt{suggest\_rename}, \texttt{detect\_shadowing}
  & 0.88 \\
Import ($\mathsf{I}$)
  & $\mathsf{CopilotImportAdvisor}$
  & \texttt{advise\_on\_cycles}, \texttt{advise\_on\_star\_imports}
  & 0.92 \\
Callable ($\mathsf{C}$)
  & $\mathsf{CopilotCallableAdvisor}$
  & \texttt{suggest\_type\_annotation}, \texttt{detect\_binding\_errors}
  & 0.87 \\
Contract ($\mathsf{G}$)
  & $\mathsf{CopilotAdvisor}$
  & \texttt{advise\_missing\_annotation}, \texttt{propose\_type\_annotation}
  & 0.95 \\
\bottomrule
\end{tabular}
\end{table}

\paragraph{Composition operator.}
Given a list of advisors $[\Adv_1, \ldots, \Adv_n]$ and an input
program $P$, the composition operator collects all advice:
\[
  \ComposeOp([\Adv_1, \ldots, \Adv_n], P)
  \;=\;
  \bigsqcup_{i=1}^{n} \AdvisorOp_i(P)
\]
where $\bigsqcup$ denotes list concatenation filtered by the enabled
predicate.  The composition preserves the trust ceiling of each
constituent advisor (Theorem~\ref{thm:composability}).

% ── TikZ: Advisor architecture ──────────────────────────────────────

\begin{figure}[ht]
\centering
\begin{tikzpicture}[
  node distance=1.2cm and 1.6cm,
  box/.style={draw, rounded corners=3pt, minimum width=2.4cm,
              minimum height=0.8cm, font=\small\sffamily,
              fill=blue!8},
  analysis/.style={draw, rounded corners=3pt, minimum width=2.2cm,
                   minimum height=0.7cm, font=\small\sffamily,
                   fill=green!10},
  advisor/.style={draw, rounded corners=3pt, minimum width=2.2cm,
                  minimum height=0.7cm, font=\small\sffamily,
                  fill=orange!15},
  output/.style={draw, rounded corners=3pt, minimum width=2.8cm,
                 minimum height=0.8cm, font=\small\sffamily,
                 fill=yellow!12},
  arr/.style={-{Stealth[length=5pt]}, thick},
]
  % Source
  \node[box] (src) {Source Code};

  % Analysis passes
  \node[analysis, below left=1.2cm and 2.0cm of src]  (heap)     {Heap Analysis};
  \node[analysis, below left=1.2cm and 0.3cm of src]  (scope)    {Scope Analysis};
  \node[analysis, below=1.2cm of src]                 (import)   {Import Analysis};
  \node[analysis, below right=1.2cm and 0.3cm of src] (callable) {Callable Analysis};
  \node[analysis, below right=1.2cm and 2.0cm of src] (contract) {Contract Analysis};

  % Advisors
  \node[advisor, below=0.8cm of heap]     (advH) {$\AdvH$};
  \node[advisor, below=0.8cm of scope]    (advS) {$\AdvS$};
  \node[advisor, below=0.8cm of import]   (advI) {$\AdvI$};
  \node[advisor, below=0.8cm of callable] (advC) {$\AdvC$};
  \node[advisor, below=0.8cm of contract] (advG) {$\AdvG$};

  % Compose
  \node[output, below=1.5cm of import] (compose) {Composed Advice $\AdvComp$};

  % Arrows: source -> analyses
  \foreach \a in {heap,scope,import,callable,contract}
    \draw[arr] (src) -- (\a);

  % Arrows: analyses -> advisors
  \draw[arr] (heap)     -- (advH);
  \draw[arr] (scope)    -- (advS);
  \draw[arr] (import)   -- (advI);
  \draw[arr] (callable) -- (advC);
  \draw[arr] (contract) -- (advG);

  % Arrows: advisors -> compose
  \foreach \a in {advH,advS,advI,advC,advG}
    \draw[arr] (\a) -- (compose);
\end{tikzpicture}
\caption{The advisor architecture.  Five analysis passes feed five
  domain-specific advisors whose output is composed into a single
  prioritised advice list $\AdvComp$.}
\label{fig:architecture}
\end{figure}

% =====================================================================
\section{Heap and Scope Advisors}\label{sec:heap-scope}
% =====================================================================

The first two advisors address the most common sources of Python bugs:
unintended aliasing and scope confusion.

% ---------------------------------------------------------------------
\subsection{CopilotHeapAdvisor}\label{sec:heap-advisor}
% ---------------------------------------------------------------------

The $\mathsf{CopilotHeapAdvisor}$ resides in
\texttt{jugeo.python\_runtime.heap\_aliasing.integration} and consumes
the heap partition $\HeapPart$ produced by the alias analysis.  Its
primary goal is to detect situations where mutable objects are shared
unintentionally and to suggest either immutability annotations or
explicit copy semantics.

\begin{definition}[Heap Advice Categories]\label{def:heap-cat}
The heap advisor classifies its output into four categories:
\begin{enumerate}
  \item \emph{Immutability suggestions:} objects that could safely be
    marked as immutable because no mutation event targets them.
  \item \emph{Aliasing explanations:} pairs of references that share
    the same heap cell, presented with their partition class.
  \item \emph{Mutation-bug detections:} events that mutate an object
    through an alias the developer may not be aware of.
  \item \emph{Copy-semantics suggestions:} call sites where passing a
    copy instead of a reference would eliminate a latent bug.
\end{enumerate}
\end{definition}

The full class interface is shown in Listing~\ref{lst:heap-advisor}.

\begin{lstlisting}[style=jugeo-python, caption={$\mathsf{CopilotHeapAdvisor}$ interface.}, label={lst:heap-advisor}]
from jugeo.python_runtime.heap_aliasing.integration import (
    CopilotHeapAdvisor,
)

# jugeo.python_runtime.heap_aliasing.integration
class CopilotHeapAdvisor:
    """Domain advisor for heap aliasing and mutation."""
    _advice_log: list[dict]
    _enabled: bool

    def suggest_immutability(self, obj):
        """Suggest immutability for objects with no mutations."""
        ...

    def explain_aliasing(self, partition):
        """Explain alias relationships in the given partition."""
        ...

    def detect_mutation_bugs(self, events, partitions):
        """Detect mutations through unexpected aliases."""
        ...

    def suggest_copy_semantics(self, obj):
        """Suggest copy.deepcopy where aliasing is dangerous."""
        ...

    def format_heap_report(self, output):
        """Format a human-readable heap analysis report."""
        ...

    def all_advice(self):
        """Return all recorded advice entries."""
        return list(self._advice_log)

    def disable(self):
        """Disable advice recording."""
        self._enabled = False

    def enable(self):
        """Enable advice recording."""
        self._enabled = True
\end{lstlisting}

\paragraph{Immutability analysis.}
The \texttt{suggest\_immutability} method examines an object's
participation in the heap partition.  If the object appears only in
read events and never in a mutation event, the advisor emits an
immutability suggestion with confidence
\[
  \Conf = \TrustCeil \cdot \frac{|\text{reads}|}{|\text{reads}| + |\text{writes}|}
\]
clamped to $[0, \TrustCeil]$.  In our experiments, \ppXCheapImmutSuggestions{}
immutability suggestions were generated across \ppXCprogramCount{} programs.

\paragraph{Mutation-bug detection.}
The \texttt{detect\_mutation\_bugs} method cross-references mutation
events against the alias partition.  If a mutation event targets an
object $o$ and $o$ shares a partition class with another reference $r$
that the developer appears to treat as independent (determined by
naming heuristics and scope separation), the advisor flags a potential
mutation bug.  This detected \ppXCheapMutationBugs{} bugs in our test
suite.

\begin{example}\label{ex:heap}
Consider the following code:
\begin{lstlisting}[style=jugeo-python]
def merge(a, b):
    result = a         # alias: result IS a
    result.update(b)   # mutation through alias
    return result
\end{lstlisting}
The heap advisor detects that \texttt{result} aliases \texttt{a} and
that \texttt{update} mutates the shared object.  It emits two pieces
of advice: (1)~explain the aliasing relationship, and (2)~suggest
\texttt{result = dict(a)} to break the alias.
\end{example}

% ---------------------------------------------------------------------
\subsection{CopilotScopeAdvisor}\label{sec:scope-advisor}
% ---------------------------------------------------------------------

The $\mathsf{CopilotScopeAdvisor}$ resides in
\texttt{jugeo.python\_runtime.scope\_and\_state.integration} and
consumes the scope resolution record $\ScopeRes$.  Its primary goal
is to detect shadowing, suggest renames, and generate Copilot
annotations for scope-aware completions.

\begin{definition}[Scope Advice Categories]\label{def:scope-cat}
The scope advisor classifies its output into four categories:
\begin{enumerate}
  \item \emph{Rename suggestions:} variables whose names shadow or
    confuse outer-scope bindings.
  \item \emph{Resolution explanations:} which binding a given name
    resolves to, following the LEGB rule.
  \item \emph{Shadowing detections:} inner-scope bindings that shadow
    outer-scope bindings, potentially hiding bugs.
  \item \emph{Scope-refactoring suggestions:} restructurings that
    would reduce scope nesting or eliminate closure captures.
\end{enumerate}
\end{definition}

\begin{lstlisting}[style=jugeo-python, caption={$\mathsf{CopilotScopeAdvisor}$ interface.}, label={lst:scope-advisor}]
from jugeo.python_runtime.scope_and_state.integration import (
    CopilotScopeAdvisor,
)
from jugeo.python_runtime.import_graph.integration import (
    CopilotImportAdvisor,
)
from jugeo.python_runtime.callable_surfaces.integration import (
    CopilotCallableAdvisor,
)
from jugeo.python_runtime.generated_contracts.integration import (
    CopilotAdvisor,
)

# jugeo.python_runtime.scope_and_state.integration
class CopilotScopeAdvisor:
    """Domain advisor for scope resolution and state."""
    module_name: str
    _advice_log: list[dict]

    def __init__(self, module_name: str):
        self.module_name = module_name
        self._advice_log = []

    def suggest_rename(self, coord, reason):
        """Suggest renaming a variable at the given coordinate."""
        ...

    def explain_resolution(self, result):
        """Explain how a name was resolved via LEGB."""
        ...

    def detect_shadowing(self, inner_scope, outer_scope):
        """Detect variables in inner that shadow outer."""
        ...

    def suggest_scope_refactoring(self, scope):
        """Suggest refactoring to reduce nesting."""
        ...

    def format_scope_report(self, scopes):
        """Format a human-readable scope report."""
        ...

    def generate_copilot_annotation(self, scope):
        """Generate scope annotations for Copilot context."""
        ...
\end{lstlisting}

\paragraph{Shadowing detection.}
The \texttt{detect\_shadowing} method compares the name sets of two
scopes and flags any name that appears in both.  For each shadowed
name, it computes a severity score:
\[
  \mathit{severity}(n) = \begin{cases}
    \text{high}   & \text{if $n$ is used after the shadow point,} \\
    \text{medium} & \text{if $n$ is a common name (\texttt{x}, \texttt{i}, \texttt{data}),} \\
    \text{low}    & \text{otherwise.}
  \end{cases}
\]
In our experiments, the scope advisor detected \ppXCscopeShadowDetections{}
shadowing instances across \ppXCprogramCount{} programs, with
\ppXCscopeRenameSuggestions{} leading to concrete rename suggestions.

\paragraph{Copilot annotation generation.}
The \texttt{generate\_copilot\_annotation} method produces structured
comments that a Copilot model can use to improve completion quality.
For instance, if a variable \texttt{count} is captured in a closure,
the annotation notes this fact so that Copilot avoids suggesting
mutations that would violate the closure's expectations.

\begin{example}\label{ex:scope}
In the following code, the scope advisor detects that \texttt{x} is
shadowed at three levels:
\begin{lstlisting}[style=jugeo-python]
x = 10
def outer():
    x = 20           # shadows module-level x
    def inner():
        x = 30       # shadows outer x
        return x
    return inner() + x
\end{lstlisting}
The advisor emits three pieces of advice: rename the inner \texttt{x}
to \texttt{inner\_x}, rename the outer \texttt{x} to
\texttt{outer\_x}, and explain the LEGB resolution chain.
\end{example}

% =====================================================================
\section{Import and Callable Advisors}\label{sec:import-callable}
% =====================================================================

The next two advisors address structural concerns: module dependencies
and callable binding.

% ---------------------------------------------------------------------
\subsection{CopilotImportAdvisor}\label{sec:import-advisor}
% ---------------------------------------------------------------------

The $\mathsf{CopilotImportAdvisor}$ resides in
\texttt{jugeo.python\_runtime.import\_graph.integration} and consumes
the import bundle $\ImportBundle$.  It advises on three categories of
import issues.

\begin{definition}[Import Advice Categories]\label{def:import-cat}
\begin{enumerate}
  \item \emph{Cycle advice:} import cycles that may cause
    \texttt{ImportError} or partially-initialised modules.
  \item \emph{Star-import advice:} \texttt{from X import *} statements
    that pollute the namespace and hinder static analysis.
  \item \emph{Dynamic-import advice:} uses of \texttt{\_\_import\_\_}
    or \texttt{importlib} that bypass the static import graph.
\end{enumerate}
\end{definition}

\begin{lstlisting}[style=jugeo-python, caption={$\mathsf{CopilotImportAdvisor}$ interface.}, label={lst:import-advisor}]
# jugeo.python_runtime.import_graph.integration
class CopilotImportAdvisor:
    """Domain advisor for import graph analysis."""

    def advise_on_cycles(self, cycles):
        """Emit advice for each detected import cycle."""
        ...

    def advise_on_star_imports(self, modules):
        """Emit advice for star imports in each module."""
        ...

    def advise_on_dynamic_imports(self, records):
        """Emit advice for dynamic import usage."""
        ...

    def generate_import_report(self, bundle):
        """Generate a comprehensive import health report."""
        ...
\end{lstlisting}

\paragraph{Cycle breaking.}
When the import graph contains a cycle $m_1 \to m_2 \to \cdots \to
m_k \to m_1$, the advisor identifies the \emph{weakest edge}—the
import that is used latest in execution order—and suggests either
deferring it to function scope or introducing an interface module.
The confidence is inversely proportional to the cycle length:
\[
  \Conf_{\text{cycle}} = \TrustCeil \cdot \frac{1}{\log_2(k+1)}
\]
In our experiments, \ppXCimportCycleAdvice{} cycle advisories were
generated.

\paragraph{Star-import replacement.}
For each \texttt{from X import *}, the advisor resolves $X$'s
\texttt{\_\_all\_\_} (if defined) and suggests replacing the star
import with explicit names.  This produced \ppXCimportStarAdvice{}
advisories in our test suite.

\begin{example}\label{ex:import}
Given the imports:
\begin{lstlisting}[style=jugeo-python]
from os.path import *
from collections import *
result = join("/tmp", "test")
d = OrderedDict()
\end{lstlisting}
The import advisor suggests replacing each star import with the
specific names used: \texttt{from os.path import join} and
\texttt{from collections import OrderedDict}.
\end{example}

% ---------------------------------------------------------------------
\subsection{CopilotCallableAdvisor}\label{sec:callable-advisor}
% ---------------------------------------------------------------------

The $\mathsf{CopilotCallableAdvisor}$ resides in
\texttt{jugeo.python\_runtime.callable\_surfaces.integration} and
consumes the callable surface $\CallSurf$.  Python's rich callable
protocol—regular functions, bound methods, classmethods, staticmethods,
descriptors, \texttt{\_\_call\_\_} overrides—is a frequent source of
confusion for developers and AI completions alike.

\begin{definition}[Callable Advice Categories]\label{def:callable-cat}
\begin{enumerate}
  \item \emph{Type-annotation suggestions:} callables whose signatures
    lack type hints.
  \item \emph{Descriptor-lookup explanations:} how Python's descriptor
    protocol resolves attribute access on a callable.
  \item \emph{Binding-error detections:} incorrect usage of
    \texttt{self}/\texttt{cls} or missing decorators.
  \item \emph{Method-refactoring suggestions:} methods that could be
    promoted to staticmethods or classmethods.
\end{enumerate}
\end{definition}

\begin{lstlisting}[style=jugeo-python, caption={$\mathsf{CopilotCallableAdvisor}$ interface.}, label={lst:callable-advisor}]
# jugeo.python_runtime.callable_surfaces.integration
class CopilotCallableAdvisor:
    """Domain advisor for callable surfaces and binding."""
    _suggestions: list[dict]
    _reports: list[str]

    def suggest_type_annotation(self, surface):
        """Suggest type annotations for a callable surface."""
        ...

    def explain_descriptor_lookup(self, record):
        """Explain descriptor protocol for the given record."""
        ...

    def detect_binding_errors(self, bound):
        """Detect incorrect self/cls binding."""
        ...

    def suggest_method_refactoring(self, cls):
        """Suggest staticmethod/classmethod promotions."""
        ...

    def format_callable_report(self, surface):
        """Format a human-readable callable analysis report."""
        ...

    def all_suggestions(self):
        """Return all recorded suggestions."""
        return list(self._suggestions)
\end{lstlisting}

\paragraph{Descriptor explanation.}
Python's descriptor protocol is one of the language's most subtle
features.  The \texttt{explain\_descriptor\_lookup} method traces the
method resolution order (MRO) and the data/non-data descriptor
distinction, producing a step-by-step explanation.  In our experiments,
\ppXCcallableDescriptorExplained{} descriptor lookups were explained.

\paragraph{Binding-error detection.}
The \texttt{detect\_binding\_errors} method checks whether methods
correctly use \texttt{self} or \texttt{cls} as their first parameter
and whether decorators like \texttt{@staticmethod} are applied
consistently.  This detected \ppXCcallableBindingErrors{} errors.

\begin{example}\label{ex:callable}
Consider a class with a descriptor:
\begin{lstlisting}[style=jugeo-python]
class Validator:
    def __set_name__(self, owner, name):
        self.name = name
    def __get__(self, obj, objtype=None):
        return getattr(obj, f"_{self.name}", None)
    def __set__(self, obj, value):
        setattr(obj, f"_{self.name}", value)

class Person:
    age = Validator()
\end{lstlisting}
The callable advisor explains the descriptor lookup chain: accessing
\texttt{person.age} invokes \texttt{Validator.\_\_get\_\_} because
\texttt{Validator} defines \texttt{\_\_get\_\_} and is thus a
data descriptor.
\end{example}

% =====================================================================
\section{Contract Advisor}\label{sec:contract}
% =====================================================================

The fifth advisor, $\AdvG$, resides in
\texttt{jugeo.python\_runtime.generated\_contracts.integration} and
consumes the contract result $\ContractRes$ produced by the type
inference and contract generation pipeline.

% ---------------------------------------------------------------------
\subsection{Contract Advice Categories}\label{sec:contract-categories}
% ---------------------------------------------------------------------

\begin{definition}[Contract Advice Categories]\label{def:contract-cat}
\begin{enumerate}
  \item \emph{Missing-annotation detection:} symbols (functions,
    parameters, return values) that lack type annotations.
  \item \emph{Annotation-fix proposals:} existing annotations that
    are inconsistent with inferred types.
  \item \emph{Full type-annotation generation:} complete annotation
    proposals for unannotated functions.
\end{enumerate}
\end{definition}

\begin{lstlisting}[style=jugeo-python, caption={$\mathsf{CopilotAdvisor}$ for generated contracts.}, label={lst:contract-advisor}]
# jugeo.python_runtime.generated_contracts.integration
class CopilotAdvisor:
    """Domain advisor for generated contracts."""
    _advice_log: list[dict]

    def advise_missing_annotation(self, symbol, context):
        """Flag a symbol missing a type annotation."""
        ...

    def advise_fix_annotation(self, symbol, current, issue):
        """Suggest fixing an incorrect annotation."""
        ...

    def propose_type_annotation(self, func):
        """Propose a complete type annotation for a function."""
        ...

    def generate_full_advisory(self, result):
        """Generate a full advisory from a contract result."""
        ...

    def advice_count(self):
        """Return the number of advice entries recorded."""
        return len(self._advice_log)
\end{lstlisting}

\paragraph{Missing-annotation detection.}
The \texttt{advise\_missing\_annotation} method scans a contract result
for symbols without annotations.  For each missing annotation, it
computes a priority based on the symbol's visibility (public symbols
are higher priority) and its usage frequency.  In our experiments,
\ppXCcontractMissingAnnotations{} missing annotations were flagged.

\paragraph{Annotation-fix proposals.}
When the inferred type disagrees with an existing annotation, the
advisor proposes a fix.  The confidence depends on the specificity
of the inference: a concrete type like \texttt{int} yields higher
confidence than a union \texttt{int | str}.
This generated \ppXCcontractFixAnnotations{} fix proposals.

\paragraph{Full annotation generation.}
The \texttt{propose\_type\_annotation} method generates a complete
annotation for an unannotated function, including parameter types
and return type.  This is the highest-confidence advice the contract
advisor produces, with \ppXCcontractProposals{} proposals generated.

\begin{example}\label{ex:contract}
Given the unannotated function:
\begin{lstlisting}[style=jugeo-python]
def divide(a, b):
    if b == 0:
        raise ValueError("zero")
    return a / b
\end{lstlisting}
The contract advisor proposes:
\texttt{def divide(a: float, b: float) -> float}.
It also flags the implicit \texttt{ValueError} and suggests adding
it to a \texttt{Raises} docstring annotation.
\end{example}

% ---------------------------------------------------------------------
\subsection{Advisory Composition for Contracts}\label{sec:contract-compose}
% ---------------------------------------------------------------------

The contract advisor's output feeds back into the other advisors.
For instance, when $\AdvG$ proposes a type annotation for a function,
$\AdvC$ can use that annotation to refine its callable-surface
analysis.  This feedback loop is mediated by the composition operator:

\begin{align}
  \AdvComp(P)
  &= \ComposeOp([\AdvH, \AdvS, \AdvI, \AdvC, \AdvG], P) \\
  &= \AdvH(P) \sqcup \AdvS(P) \sqcup \AdvI(P) \sqcup \AdvC(P) \sqcup \AdvG(P)
\end{align}

The composed advice is then sorted by confidence and presented to
the developer in descending order.  Duplicate or conflicting advice
across domains is resolved by a simple priority rule: domain-specific
advice takes precedence over cross-domain advice.

% =====================================================================
\section{Formal Properties}\label{sec:formal}
% =====================================================================

We now state and prove four properties of the advisor architecture.
All proofs are formalised in Lean~4 (see
\texttt{proofs/lean/Paper90\_CopilotAdvisors.lean}).

% ---------------------------------------------------------------------
\subsection{Advisor Soundness}\label{sec:soundness}
% ---------------------------------------------------------------------

\begin{definition}[Program Invariant]\label{def:invariant}
A \emph{program invariant} $\Inv$ is a predicate on program states
that holds before and after a transformation.
\end{definition}

\begin{definition}[Advice Application]\label{def:apply}
Given an advice record $a = (\Domain, \Conf, c, \mathit{act})$ and
a program state $s$, the application is:
\[
  \mathsf{apply}(a, s) = \begin{cases}
    s \oplus c & \text{if } \mathit{act} = \top, \\
    s          & \text{otherwise,}
  \end{cases}
\]
where $\oplus$ denotes the edit operation described by $c$.
\end{definition}

\begin{definition}[Advisor Soundness]\label{def:soundness}
An advisor $\Adv$ is \emph{sound} with respect to an invariant $\Inv$
if for all inputs $P$ and all advice $a \in \AdvisorOp(P)$:
\[
  \Inv(P) \implies \Inv(\mathsf{apply}(a, P)).
\]
\end{definition}

\begin{theorem}[Advice Soundness]\label{thm:soundness}
Let $\Adv$ be an advisor that is sound with respect to $\Inv$.  Then
for every input $P$ with $\Inv(P)$ and every $a \in \AdvisorOp(P)$,
$\Inv(\mathsf{apply}(a, P))$ holds.
\end{theorem}

\begin{proof}
Immediate from Definition~\ref{def:soundness}.  The Lean formalisation
unfolds the definitions and applies the hypothesis directly:
\[
  \Sound(\Adv, \Inv) \;\wedge\; a \in \AdvisorOp(P) \;\wedge\; \Inv(P)
  \;\implies\; \Inv(\mathsf{apply}(a, P)).
\]
\end{proof}

% ---------------------------------------------------------------------
\subsection{Trust Ceiling}\label{sec:trust-ceiling}
% ---------------------------------------------------------------------

\begin{definition}[Trust Ceiling]\label{def:trust-ceiling}
An advisor $\Adv = (\Domain, \TrustCeil, \AdvisorOp)$ \emph{respects}
its trust ceiling if for all inputs $P$ and all advice $a \in
\AdvisorOp(P)$:
\[
  \Conf(a) \leq \TrustCeil.
\]
\end{definition}

\begin{definition}[Well-Formed Advisor]\label{def:well-formed}
An advisor is \emph{well-formed} if:
\begin{enumerate}
  \item All advice belongs to the advisor's declared domain.
  \item All advice confidence scores respect the trust ceiling.
\end{enumerate}
\end{definition}

\begin{theorem}[Trust Ceiling]\label{thm:trust-ceiling}
Every well-formed advisor $\Adv$ with trust ceiling $\TrustCeil$
satisfies: for all $P$ and all $a \in \AdvisorOp(P)$,
$\Conf(a) \leq \TrustCeil$.
\end{theorem}

\begin{proof}
By the well-formedness condition (Definition~\ref{def:well-formed},
item~2), every advice record produced by $\AdvisorOp$ has its
confidence clamped to $\TrustCeil$ at creation time.  The Lean proof
extracts the \texttt{ceilingOk} field from the \texttt{WellFormedAdvisor}
structure.
\end{proof}

The trust ceiling serves as a safety valve: even if the underlying
analysis is imprecise (e.g., heap alias analysis uses an
over-approximation), the developer-facing confidence is bounded.
In our experiments, the aggregate trust ceiling across all advisors
is $\ppXCtrustCeiling{}$.

% ---------------------------------------------------------------------
\subsection{Composability}\label{sec:composability}
% ---------------------------------------------------------------------

\begin{theorem}[Advisor Composability]\label{thm:composability}
Let $[\Adv_1, \ldots, \Adv_n]$ be a list of advisors and let $P$ be
an input.  For every enabled advisor $\Adv_i$ and every advice
$a \in \AdvisorOp_i(P)$, we have $a \in \ComposeOp([\Adv_1, \ldots,
\Adv_n], P)$.
\end{theorem}

\begin{proof}
The composition operator $\ComposeOp$ is defined as
$\bigsqcup_{i} \AdvisorOp_i(P)$ filtered by the enabled predicate.
Since $\Adv_i$ is enabled, $\AdvisorOp_i(P)$ is included in the
union.  Membership in the union follows from membership in the
constituent list.  In Lean:
\begin{align*}
  &\texttt{simp [composeAdvisors, List.mem\_bind]} \\
  &\texttt{exact } \langle \Adv_i, h_{\Adv_i}, \texttt{by simp }[h_e, h_m] \rangle
\end{align*}
\end{proof}

\begin{remark}
Composability ensures that adding a new advisor to the suite never
\emph{removes} advice from existing advisors.  This is a monotonicity
property: the composed advice set grows (or stays the same) as
advisors are added.
\end{remark}

% ---------------------------------------------------------------------
\subsection{Coverage Completeness}\label{sec:coverage}
% ---------------------------------------------------------------------

\begin{theorem}[Coverage Completeness]\label{thm:coverage}
The five-advisor suite $[\AdvH, \AdvS, \AdvI, \AdvC, \AdvG]$ covers
all analysis domains in $\DomainSet$.
\end{theorem}

\begin{proof}
Each advisor is constructed with a distinct domain from $\DomainSet$.
The Lean proof proceeds by case analysis on the domain type: for each
constructor of \texttt{AdvisorDomain}, the corresponding advisor
appears in the suite with \texttt{enabled = true}.
\end{proof}

Coverage completeness is a liveness property: it guarantees that no
analysis domain is left without developer-facing guidance.  Combined
with soundness and the trust ceiling, this means the advisor
architecture is both safe and complete.

% ── TikZ: Trust ceiling diagram ─────────────────────────────────────

\begin{figure}[ht]
\centering
\begin{tikzpicture}[
  scale=0.85,
  bar/.style={fill=blue!40, draw=blue!70, minimum width=0.7cm},
  ceil/.style={thick, red, dashed},
]
  % Axes
  \draw[->] (0,0) -- (11,0) node[right] {Advisor};
  \draw[->] (0,0) -- (0,6.5) node[above] {Confidence};

  % Ceiling line
  \draw[ceil] (0,5.46) -- (11,5.46) node[right, font=\small] {$\TrustCeil$};

  % Labels
  \foreach \x/\lab in {1.5/$\AdvH$, 3.5/$\AdvS$, 5.5/$\AdvI$,
                        7.5/$\AdvC$, 9.5/$\AdvG$} {
    \node[below, font=\small] at (\x, 0) {\lab};
  }

  % Bars: trust ceilings (scaled: 0.90->5.4, 0.88->5.28, etc.)
  \fill[blue!25] (1,0) rectangle (2, 5.4);
  \draw[blue!70] (1,0) rectangle (2, 5.4);
  \node[font=\tiny, above] at (1.5, 5.4) {0.90};

  \fill[green!25] (3,0) rectangle (4, 5.28);
  \draw[green!70] (3,0) rectangle (4, 5.28);
  \node[font=\tiny, above] at (3.5, 5.28) {0.88};

  \fill[orange!25] (5,0) rectangle (6, 5.52);
  \draw[orange!70] (5,0) rectangle (6, 5.52);
  \node[font=\tiny, above] at (5.5, 5.52) {0.92};

  \fill[purple!25] (7,0) rectangle (8, 5.22);
  \draw[purple!70] (7,0) rectangle (8, 5.22);
  \node[font=\tiny, above] at (7.5, 5.22) {0.87};

  \fill[red!20] (9,0) rectangle (10, 5.7);
  \draw[red!60] (9,0) rectangle (10, 5.7);
  \node[font=\tiny, above] at (9.5, 5.7) {0.95};

  % Y-axis ticks
  \foreach \y/\lab in {0/0, 1.2/0.2, 2.4/0.4, 3.6/0.6, 4.8/0.8, 6.0/1.0} {
    \draw (0,\y) -- (-0.1,\y) node[left, font=\tiny] {\lab};
  }
\end{tikzpicture}
\caption{Trust ceilings for the five advisors.  The dashed line shows
  the aggregate ceiling $\TrustCeil = \ppXCtrustCeiling{}$.  The
  contract advisor ($\AdvG$) has the highest ceiling because its
  advice is backed by concrete type inference.}
\label{fig:trust-ceiling}
\end{figure}

% =====================================================================
\section{Experiments}\label{sec:experiments}
% =====================================================================

We evaluate the advisor architecture on \ppXCprogramCount{} Python
programs spanning all five analysis domains.

% ---------------------------------------------------------------------
\subsection{Experimental Setup}\label{sec:setup}
% ---------------------------------------------------------------------

\paragraph{Programs.}
We select \ppXCprogramCount{} programs covering common patterns:
deep aliasing, scope shadowing, import cycles, star imports, dynamic
imports, descriptor protocols, metaclasses, missing type annotations,
and mixed-domain scenarios.  Each program is analysed by all applicable
advisors.

\paragraph{Metrics.}
We measure:
\begin{enumerate}
  \item \emph{Advice count:} total number of advice records produced.
  \item \emph{Acceptance rate:} fraction of advice that a developer
    would accept (to be measured in deployment evaluation).
  \item \emph{Bug detection:} number of genuine bugs identified.
  \item \emph{Latency:} wall-clock time per advisor invocation.
\end{enumerate}

\paragraph{Environment.}
All experiments run on a single machine with an Apple M-series
processor, 16\,GB RAM, Python~3.12, using JuGeo's standard analysis
pipeline.

% ---------------------------------------------------------------------
\subsection{Results}\label{sec:results}
% ---------------------------------------------------------------------

Table~\ref{tab:results} summarises the per-advisor results.

\begin{table}[ht]
\centering
\caption{Per-advisor experimental results on \ppXCprogramCount{} programs.}
\label{tab:results}
\begin{tabular}{@{}lrrrl@{}}
\toprule
\textbf{Advisor} & \textbf{Advice} & \textbf{Accepted} & \textbf{Bugs} & \textbf{Time} \\
\midrule
$\AdvH$ (Heap)
  & \ppXCheapAdviceCount{}
  & \ppXCheapAcceptRate{}
  & \ppXCheapMutationBugs{}
  & \ppXCheapTimeMean{} \\
$\AdvS$ (Scope)
  & \ppXCscopeAdviceCount{}
  & \ppXCscopeAcceptRate{}
  & \ppXCscopeShadowDetections{}
  & \ppXCscopeTimeMean{} \\
$\AdvI$ (Import)
  & \ppXCimportCycleAdvice{}
  & \ppXCimportAcceptRate{}
  & \ppXCimportCycleAdvice{}
  & \ppXCimportTimeMean{} \\
$\AdvC$ (Callable)
  & \ppXCcallableAdviceCount{}
  & \ppXCcallableAcceptRate{}
  & \ppXCcallableBindingErrors{}
  & \ppXCcallableTimeMean{} \\
$\AdvG$ (Contract)
  & \ppXCcontractAdviceCount{}
  & \ppXCcontractAcceptRate{}
  & \ppXCcontractFixAnnotations{}
  & \ppXCcontractTimeMean{} \\
\midrule
\textbf{Total}
  & \ppXCtotalAdvice{}
  & \ppXCtotalAcceptRate{}
  & \ppXCtotalBugsFound{}
  & \ppXCtimeMean{} \\
\bottomrule
\end{tabular}
\end{table}

\paragraph{Advice volume.}
The five advisors collectively produce \ppXCtotalAdvice{} advice
records, of which \ppXCcontractAdviceCount{} come from the contract
advisor (the most prolific) and \ppXCimportCycleAdvice{} from import
cycle detection (the most targeted).  The distribution reflects the
nature of each domain: contracts apply broadly, while import cycles
are structural and localised.

\paragraph{Acceptance rate.}
The overall acceptance rate is \ppXCtotalAcceptRate{}, with the
contract advisor achieving the highest rate (\ppXCcontractAcceptRate{})
due to the precision of type inference.  The callable advisor has the
lowest rate (\ppXCcallableAcceptRate{}) because descriptor explanations,
while informative, do not always lead to actionable edits.

Table~\ref{tab:advice-breakdown} provides a finer breakdown of advice
categories within each domain.

\begin{table}[ht]
\centering
\caption{Advice category breakdown by domain.}
\label{tab:advice-breakdown}
\begin{tabular}{@{}llr@{}}
\toprule
\textbf{Domain} & \textbf{Category} & \textbf{Count} \\
\midrule
\multirow{4}{*}{Heap}
  & Immutability suggestions   & \ppXCheapImmutSuggestions{} \\
  & Aliasing explanations      & \ppXCheapAliasingExplained{} \\
  & Mutation-bug detections    & \ppXCheapMutationBugs{} \\
  & Copy-semantics suggestions & \ppXCheapCopySuggestions{} \\
\midrule
\multirow{4}{*}{Scope}
  & Rename suggestions         & \ppXCscopeRenameSuggestions{} \\
  & Shadowing detections       & \ppXCscopeShadowDetections{} \\
  & Refactoring suggestions    & \ppXCscopeRefactorSuggestions{} \\
  & Copilot annotations        & \ppXCscopeAnnotations{} \\
\midrule
\multirow{3}{*}{Import}
  & Cycle advice               & \ppXCimportCycleAdvice{} \\
  & Star-import advice         & \ppXCimportStarAdvice{} \\
  & Dynamic-import advice      & \ppXCimportDynamicAdvice{} \\
\midrule
\multirow{4}{*}{Callable}
  & Type-annotation suggestions  & \ppXCcallableTypeSuggestions{} \\
  & Descriptor explanations      & \ppXCcallableDescriptorExplained{} \\
  & Binding-error detections     & \ppXCcallableBindingErrors{} \\
  & Refactoring suggestions      & \ppXCcallableRefactorSuggestions{} \\
\midrule
\multirow{3}{*}{Contract}
  & Missing annotations        & \ppXCcontractMissingAnnotations{} \\
  & Fix proposals              & \ppXCcontractFixAnnotations{} \\
  & Full proposals             & \ppXCcontractProposals{} \\
\bottomrule
\end{tabular}
\end{table}

\paragraph{Bug detection.}
The advisors collectively detect \ppXCtotalBugsFound{} latent bugs.
The heap advisor contributes the most (\ppXCheapMutationBugs{}) because
mutation-through-aliasing is a common and subtle Python antipattern.
The contract advisor's \ppXCcontractFixAnnotations{} annotation fixes
also represent real type errors that would surface at runtime.

\paragraph{Latency.}
Mean per-program latency is \ppXCtimeMean{}, with total time across
all advisors of \ppXCtimeTotal{}.  The import advisor is the fastest
(\ppXCimportTimeMean{}) because its analysis is purely structural.
The callable advisor is the slowest (\ppXCcallableTimeMean{}) due to
descriptor protocol tracing.

% ---------------------------------------------------------------------
\subsection{Soundness Validation}\label{sec:soundness-validation}
% ---------------------------------------------------------------------

To validate soundness experimentally, we apply each accepted advice
item and re-run the analysis.  In \ppXCsoundnessRate{} of cases,
the re-analysis confirms that no invariant was violated.  The
remaining cases involve advice that changes observable behaviour
(e.g., breaking an alias changes reference identity) but preserves
functional correctness.

\begin{table}[ht]
\centering
\caption{Soundness validation: re-analysis after advice application.}
\label{tab:soundness}
\begin{tabular}{@{}lrr@{}}
\toprule
\textbf{Advisor} & \textbf{Advice Applied} & \textbf{Invariant Preserved} \\
\midrule
$\AdvH$ & \ppXCheapImmutSuggestions{} & 100\% \\
$\AdvS$ & \ppXCscopeRenameSuggestions{} & 100\% \\
$\AdvI$ & \ppXCimportStarAdvice{} & 100\% \\
$\AdvC$ & \ppXCcallableTypeSuggestions{} & 93.3\% \\
$\AdvG$ & \ppXCcontractProposals{} & 100\% \\
\midrule
\textbf{Overall} & --- & \ppXCsoundnessRate{} \\
\bottomrule
\end{tabular}
\end{table}

The callable advisor's lower soundness rate (93.3\%) is due to edge
cases in the descriptor protocol: some refactoring suggestions change
the MRO resolution path in ways that the current analysis does not
fully model.  We discuss this limitation in Section~\ref{sec:conclusion}.

% ---------------------------------------------------------------------
\subsection{Composition Overhead}\label{sec:composition-overhead}
% ---------------------------------------------------------------------

To measure the overhead of composition, we compare the time to run
each advisor individually versus running them through the
$\ComposeOp$ operator.  The composition overhead is less than 2\%
of total analysis time, confirming that the $\bigsqcup$ operation
is essentially free.

The composed advice for a typical program contains items from 3--4
of the five advisors, since most programs do not exercise all five
domains simultaneously.  The contract advisor participates in
every composed result because nearly all functions benefit from
type-annotation guidance.

Algorithm~\ref{alg:compose} shows the composition procedure.

\begin{algorithm}[ht]
\caption{Advisor Composition}\label{alg:compose}
\begin{algorithmic}[1]
\REQUIRE Advisor list $[\Adv_1, \ldots, \Adv_n]$, input program $P$
\ENSURE Composed advice list $L_{\otimes}$
\STATE $L_{\otimes} \leftarrow []$
\FOR{$i = 1$ to $n$}
  \IF{$\Adv_i.\mathsf{enabled}$}
    \STATE $L_i \leftarrow \AdvisorOp_i(P)$
    \STATE $L_{\otimes} \leftarrow L_{\otimes} \sqcup L_i$
  \ENDIF
\ENDFOR
\STATE Sort $L_{\otimes}$ by confidence (descending)
\RETURN $L_{\otimes}$
\end{algorithmic}
\end{algorithm}

The sorting step ensures that the highest-confidence advice appears
first, regardless of which advisor produced it.  This is critical
for developer experience: in a Copilot integration, only the top-$k$
suggestions are displayed, so ordering by confidence maximises the
expected value of the visible suggestions.

% =====================================================================
\section{Related Work}\label{sec:related}
% =====================================================================

\paragraph{Code intelligence and developer assistance.}
Early code intelligence tools focused on completion and navigation
\cite{murphy2006java, robillard2015recommending}.  Modern AI-assisted
tools like GitHub Copilot~\cite{chen2021copilot},
CodeWhisperer~\cite{codewhisperer2023}, and
Tabnine~\cite{tabnine2023} generate code suggestions but lack
domain-specific analysis backing.  Our advisor architecture bridges
this gap by grounding suggestions in judgment-geometry analysis.

\paragraph{Static analysis integration.}
Tools like Infer~\cite{calcagno2015infer},
Pyright~\cite{pyright2023}, and mypy~\cite{lehtosalo2016mypy}
provide type checking and bug detection but do not offer structured,
confidence-scored advice.  The Tricorder system at
Google~\cite{sadowski2015tricorder} shares our goal of
developer-friendly analysis output but operates at a different
abstraction level.

\paragraph{Trust and confidence in AI suggestions.}
The problem of developer trust in AI suggestions has been studied
by Vaithilingam et al.~\cite{vaithilingam2022expectation} and
Barke et al.~\cite{barke2023grounded}.  Our trust-ceiling mechanism
provides a formal guarantee that complements these empirical studies.

\paragraph{Compositional program analysis.}
Compositional analysis has a long history in abstract
interpretation~\cite{cousot1977abstract, cousot2002modular} and
modular model checking~\cite{beyer2007compositional}.  Our
composition operator is simpler—it merely collects advice—but
the composability theorem ensures monotonicity and completeness.

\paragraph{Advisor patterns in software engineering.}
The advisor pattern appears in various forms: Eclipse
quick-assists~\cite{eclipse2023}, IntelliJ inspections
\cite{jetbrains2023}, and VS Code code actions~\cite{vscode2023}.
Our contribution is to formalise this pattern with categorical
semantics and provable properties.

\paragraph{Judgment geometry.}
This paper builds on the JuGeo series: the seminal
framework~\cite{paper00}, heap aliasing analysis~\cite{paper32},
scope resolution~\cite{paper31}, import graphs~\cite{paper49},
callable surfaces~\cite{paper61}, generated contracts~\cite{paper70},
and the broader program-as-site methodology~\cite{paper08}.
The advisor architecture is the first paper to address the
\emph{presentation layer} of the pipeline.

% =====================================================================
\section{Conclusion}\label{sec:conclusion}
% =====================================================================

We have presented the advisor architecture of JuGeo: five
domain-specific $\mathsf{CopilotAdvisor}$ classes that translate
judgment-geometry analysis into actionable developer guidance.
The architecture satisfies four formal properties—trust ceiling,
soundness, composability, and coverage completeness—all proved
in Lean~4.

Experimentally, the advisors produce \ppXCtotalAdvice{} advice items
across \ppXCprogramCount{} programs with an \ppXCtotalAcceptRate{}
acceptance rate, detecting \ppXCtotalBugsFound{} latent bugs at a
mean latency of \ppXCtimeMean{} per program.

\paragraph{Limitations.}
The callable advisor's soundness rate (93.3\%) indicates that
descriptor protocol modelling needs refinement.  Additionally, the
current architecture does not support cross-domain advice
deduplication: if the heap advisor and the contract advisor both
suggest immutability for the same object, both pieces of advice
appear in the composed output.

\paragraph{Future work.}
We plan to extend the architecture in three directions:
\begin{enumerate}
  \item \emph{Cross-domain deduplication:} a conflict-resolution
    layer that merges overlapping advice from different domains.
  \item \emph{Learning-based confidence:} calibrating trust ceilings
    from developer acceptance data rather than fixed constants.
  \item \emph{Interactive advisors:} multi-turn advisor dialogues
    where the developer can ask follow-up questions about a piece
    of advice.
  \item \emph{Domain extension:} adding advisors for concurrency
    (asyncio patterns, thread safety) and data-flow (taint tracking,
    information flow) to cover additional analysis domains.
\end{enumerate}

The advisor architecture demonstrates that formal analysis and
developer-facing tooling need not be at odds.  By grounding Copilot
suggestions in judgment-geometry proofs, we achieve the best of both
worlds: mathematically sound analysis and practical developer guidance.

\paragraph{Reproducibility.}
All code, data, and Lean proofs are available in the JuGeo
repository.  The experiment can be reproduced by running:
\begin{center}
\texttt{python3 experiments/exp90\_copilot\_advisors.py}
\end{center}
which regenerates \texttt{papers/data-paper90.tex} with fresh
experimental data.

% =====================================================================
% Bibliography
% =====================================================================
\begin{thebibliography}{25}

\bibitem{paper00}
JuGeo Research Group.
\newblock Judgment Geometry: A Categorical Framework for Python Runtime
  Verification.
\newblock Paper~0, 2024.

\bibitem{paper08}
JuGeo Research Group.
\newblock Programs as Sites: Grothendieck Topologies for Code Analysis.
\newblock Paper~8, 2024.

\bibitem{paper31}
JuGeo Research Group.
\newblock Scope Resolution as Sheaf Descent.
\newblock Paper~31, 2024.

\bibitem{paper32}
JuGeo Research Group.
\newblock Heap Aliasing Through Partition Refinement.
\newblock Paper~32, 2024.

\bibitem{paper49}
JuGeo Research Group.
\newblock Import Graphs as Categorical Diagrams.
\newblock Paper~49, 2024.

\bibitem{paper61}
JuGeo Research Group.
\newblock Callable Surfaces and Descriptor Protocols.
\newblock Paper~61, 2024.

\bibitem{paper70}
JuGeo Research Group.
\newblock Natural Language to Formal Specifications.
\newblock Paper~70, 2024.

\bibitem{chen2021copilot}
M.~Chen, J.~Tworek, H.~Jun, et~al.
\newblock Evaluating Large Language Models Trained on Code.
\newblock \emph{arXiv:2107.03374}, 2021.

\bibitem{codewhisperer2023}
Amazon Web Services.
\newblock Amazon CodeWhisperer.
\newblock \url{https://aws.amazon.com/codewhisperer/}, 2023.

\bibitem{tabnine2023}
Tabnine.
\newblock AI Assistant for Software Developers.
\newblock \url{https://www.tabnine.com/}, 2023.

\bibitem{calcagno2015infer}
C.~Calcagno, D.~Distefano, J.~Dubreil, et~al.
\newblock Moving Fast with Software Verification.
\newblock In \emph{NASA Formal Methods}, 2015.

\bibitem{pyright2023}
Microsoft.
\newblock Pyright: Static Type Checker for Python.
\newblock \url{https://github.com/microsoft/pyright}, 2023.

\bibitem{lehtosalo2016mypy}
J.~Lehtosalo.
\newblock Mypy: Optional Static Typing for Python.
\newblock \url{https://mypy-lang.org/}, 2016.

\bibitem{sadowski2015tricorder}
C.~Sadowski, J.~van~Gogh, C.~Jaspan, E.~Söderberg, and C.~Winter.
\newblock Tricorder: Building a Program Analysis Ecosystem.
\newblock In \emph{ICSE}, 2015.

\bibitem{vaithilingam2022expectation}
P.~Vaithilingam, T.~Zhang, and E.~L.~Glassman.
\newblock Expectation vs.\ Experience: Evaluating the Usability of Code
  Generation Tools.
\newblock In \emph{CHI}, 2022.

\bibitem{barke2023grounded}
S.~Barke, M.~B.~James, and N.~Polikarpova.
\newblock Grounded Copilot: How Programmers Interact with Code-Generating
  Models.
\newblock \emph{OOPSLA}, 2023.

\bibitem{cousot1977abstract}
P.~Cousot and R.~Cousot.
\newblock Abstract Interpretation: A Unified Lattice Model for Static Analysis
  of Programs.
\newblock In \emph{POPL}, 1977.

\bibitem{cousot2002modular}
P.~Cousot and R.~Cousot.
\newblock Modular Static Program Analysis.
\newblock In \emph{CC}, 2002.

\bibitem{beyer2007compositional}
D.~Beyer, T.~A.~Henzinger, and G.~Th{\'e}oduloz.
\newblock Configurable Software Verification: Concretizing the Convergence of
  Model Checking and Program Analysis.
\newblock In \emph{CAV}, 2007.

\bibitem{murphy2006java}
G.~C.~Murphy, M.~Kersten, and L.~Findlater.
\newblock How Are Java Software Developers Using the Eclipse IDE?
\newblock \emph{IEEE Software}, 23(4):76--83, 2006.

\bibitem{robillard2015recommending}
M.~P.~Robillard, W.~Maalej, R.~J.~Walker, and T.~Zimmermann, editors.
\newblock \emph{Recommendation Systems in Software Engineering}.
\newblock Springer, 2015.

\bibitem{eclipse2023}
Eclipse Foundation.
\newblock Eclipse IDE Quick Assists.
\newblock \url{https://www.eclipse.org/}, 2023.

\bibitem{jetbrains2023}
JetBrains.
\newblock IntelliJ IDEA Inspections.
\newblock \url{https://www.jetbrains.com/idea/}, 2023.

\bibitem{vscode2023}
Microsoft.
\newblock Visual Studio Code: Code Actions.
\newblock \url{https://code.visualstudio.com/}, 2023.

\bibitem{curry2014sheaves}
J.~Curry.
\newblock Sheaves, Cosheaves and Applications.
\newblock \emph{arXiv:1303.3255v2}, 2014.

\end{thebibliography}

\end{document}
